\PassOptionsToPackage{noend}{algorithm2e}
\documentclass[pmlr]{jmlr}% new name PMLR (Proceedings of Machine Learning Research)

 % The following packages will be automatically loaded:
 % amsmath, amssymb, natbib, graphicx, url, algorithm2e

 %\usepackage{rotating}% for sideways figures and tables
\usepackage{longtable}% for long tables

 % The booktabs package is used by this sample document
 % (it provides \toprule, \midrule and \bottomrule).
 % Remove the next line if you don't require it.
\usepackage{booktabs}
 % The siunitx package is used by this sample document
 % to align numbers in a column by their decimal point.
 % Remove the next line if you don't require it.
\usepackage[load-configurations=version-1]{siunitx} % newer version
 %\usepackage{siunitx}

 % The following command is just for this sample document:
\newcommand{\cs}[1]{\texttt{\char`\\#1}}

 % Define an unnumbered theorem just for this sample document:
\theorembodyfont{\upshape}
\theoremheaderfont{\scshape}
\theorempostheader{:}
\theoremsep{\newline}
\newtheorem*{note}{Note}

 % change the arguments, as appropriate, in the following:
\jmlrvolume{1}
\jmlryear{2010}
\jmlrworkshop{Workshop Title}


% Han's packages
\usepackage{tikz}
\usetikzlibrary{arrows.meta, positioning, matrix, tikzmark}
\usepackage{subcaption}
\newcommand{\refs}{%
    \textcolor{red}{\textbf{[NEED REF]}}%
    }
\usepackage[table]{xcolor}




\title{Association Transformations in Multi-tier Structures}

%\title[Short Title]{Full Title of Article\titlebreak This Title Has
%A Line Break\titletag{\thanks{sample footnote}}}
%
 % Use \Name{Author Name} to specify the name.

 % Spaces are used to separate forenames from the surname so that
 % the surnames can be picked up for the page header and copyright footer.
 
 % If the surname contains spaces, enclose the surname
 % in braces, e.g. \Name{John {Smith Jones}} similarly
 % if the name has a "von" part, e.g \Name{Jane {de Winter}}.
 % If the first letter in the forenames is a diacritic
 % enclose the diacritic in braces, e.g. \Name{{\'E}louise Smith}

 % *** Make sure there's no spurious space before \nametag ***

 % Two authors with the same address
%  \author{\Name{Author Name1\nametag{\thanks{with a note}}} \Email{abc@sample.com}\and
%   \Name{Author Name2} \Email{xyz@sample.com}\\
%   \addr Address}

 % Three or more authors with the same address:
 % \author{\Name{Author Name1} \Email{an1@sample.com}\\
 %  \Name{Author Name2} \Email{an2@sample.com}\\
 %  \Name{Author Name3} \Email{an3@sample.com}\\
 %  \Name{Author Name4} \Email{an4@sample.com}\\
 %  \Name{Author Name5} \Email{an5@sample.com}\\
 %  \Name{Author Name6} \Email{an6@sample.com}\\
 %  \Name{Author Name7} \Email{an7@sample.com}\\
 %  \Name{Author Name8} \Email{an8@sample.com}\\
 %  \Name{Author Name9} \Email{an9@sample.com}\\
 %  \Name{Author Name10} \Email{an10@sample.com}\\
 %  \Name{Author Name11} \Email{an11@sample.com}\\
 %  \Name{Author Name12} \Email{an12@sample.com}\\
 %  \Name{Author Name13} \Email{an13@sample.com}\\
 %  \Name{Author Name14} \Email{an14@sample.com}\\
 %  \addr Address}


 % Authors with different addresses:
 % \author{\Name{Author Name1} \Email{abc@sample.com}\\
 % \addr Address 1
 % \AND
 % \Name{Author Name2} \Email{xyz@sample.com}\\
 % \addr Address 2
 %}

%\editor{Editor's name}
 % \editors{List of editors' names}

\begin{document}

\maketitle

%\begin{abstract}
%This is the abstract for this article.
%\end{abstract}
%\begin{keywords}
%List of keywords
%\end{keywords}


\section{Introduction}
This paper presents a method for learning transformations over multi-tier 
representations, with a particular focus on the modification of cross-tier 
relations (i.e., associations between tiers), which is widely used in 
linguistics. 
We consider a specific but linguistically important type of transformations where
the nodes remain unchanged from input to output, while only the association relation is
modified subject to non-crossing constraints. We formalize these 
representations as temporally-ordered tiers with a binary association 
relation, and the learning problem is
therefore framed as inferring the which association relations are
preserved, deleted, or newly introduced from observed input-output pairs.
\begin{figure}[ht!]
	\centering
	
	\begin{tikzpicture}[
		node/.style={inner sep=1.5pt, minimum size=16pt},
		assoc/.style={thick},
		scale=0.9
		]
		
		% ================= INPUT =================
		
		% top tier
		\node[node] (a1) at (0,1.5) {$a_1$};
		\node[node] (a2) at (1.2,1.5) {$a_2$};
		\node[node] (a3) at (2.4,1.5) {$a_3$};
		
		% bottom tier
		\node[node] (b1) at (0,0) {$b_1$};
		\node[node] (b2) at (1.2,0) {$b_2$};
		\node[node] (b3) at (2.4,0) {$b_3$};
		
		% associations
		\draw[assoc] (a1) -- (b1);
		\draw[assoc] (a2) -- (b2);
		\draw[assoc] (a3) -- (b3);
		
		% subcaption
		\node at (1.2,-0.8) {Input};
		
		% ================= ARROW =================
		
		\draw[->, thick] (3.5,0.75) -- (4.7,0.75);
		
		% ================= OUTPUT =================
		
		% top tier
		\node[node] (c1) at (5.8,1.5) {$a_1$};
		\node[node] (c2) at (7.0,1.5) {$a_2$};
		\node[node] (c3) at (8.2,1.5) {$a_3$};
		
		% bottom tier
		\node[node] (d2) at (7.0,0) {$b_2$};
		\node[node] (d3) at (8.2,0) {$b_3$};
		
		% associations
		\draw[assoc] (c1) -- (d2);
		\draw[assoc] (c2) -- (d3);
		\draw[assoc] (c3) -- (d3);
		
		% subcaption
		\node at (7.0,-0.8) {Output};
		
	\end{tikzpicture}
	
	\caption{
		Learning association transformations between tiers.
	}
\end{figure}

Such multi-tiered structures have been widely used in linguistic 
theory, such as tone-syllable associations in autosegmental 
phonology \citep{goldsmith1976autosegmental}
prosodic or syllable structure \citep{hayes1980metrical,
	sagey1986representation}. This insight implies that 
Linguistic information of different aspects are organized across multiple 
temporally ordered tiers. Associations between tiers encode which 
elements are linked at a given point in time, thereby capturing 
relationships across different dimensions of the structure. We
refer this type of graph as \textit{autosegmental graphs} (ASGs),
which can be viewed as bipartite graphs, where 
nodes with shared properties are organized into the same tiers. 
Since each tier is linearly ordered, associations are subject to 
non-crossing constraint \citep[NCC,][]{bird1994one,goldsmith1976autosegmental}.
The core idea behind autosegmental graphs is that input-output 
mappings are in fact more naturally characterized by rearranging 
association lines across tiers rather than changing elements themselves.
A large body of theoretical work on African tonal processes has demonstrated the 
advantages of autosegmental graphs \citep{hyman2009representation}. 
%clements2010autosegmental, oddenintroducing2013}.  
Computationally, 
using multi-tiered structures has been shown to reduce the complexity of 
tonotactic patterns \citep{jardineComputationallyTone2016}.


Previous work has modeled these structures using finite-state 
approaches \citep{wiebe1992modelling,bird1994one,yli2013finite,
	yli2015three}. These approaches typically convert structured 
representations into strings or synchronized tapes to leverage 
finite-state machinery, or rely on hard-coded mechanisms to 
represent different types of associations.

In contrast, we propose an approach that learns association 
transformations between tiers directly over autosegmental 
graphs using model theory and formal logic. As a preliminary 
study, we begin with an idealized assumption: tonal processes 
involve no changes to segmental content, but only the 
rearrangement of association lines between the underlying 
representation (UR) and the surface representation (SR). 
In addition, we assume that only a single process constitutes 
the learning target. 

We show that, by adopting autosegmental graphs within a 
model-theoretic framework of formal transductions, it is 
possible to construct a partial order over structures and 
apply a model-theoretic learner such as the Bottom-Up Factor 
Inference Algorithm \citep{chandleeBufia2019,rawski2021structure}. 
The target output of the learner is a formula that preserves 
unchanged associations, deletes associations that must be delinked, and 
introduces newly added associations. Learning proceeds through iterative 
hypothesis updates while traversing a partially ordered space of logically possible structures.

Furthermore, transformations that appear non-deterministic when modeled 
over strings can instead be expressed deterministically within this 
relational framework. To our knowledge, the problem of learning local 
transformations of association relations in multi-tier structures has not previously been addressed in this form.


\section{Preliminaries}
\subsection{Autosegmental graphs}

\begin{definition}[Autosegmental Graph]
	An autosegmental graph (ASG) is a labeled bipartite graph
	with multiple linearly-ordered tiers. Equivalently, it can be defined
	model-theoretically as a relational structure
	$\mathcal{A} = \langle D, \mathcal{T}, \alpha, \lhd \rangle$
	where $D$ is a finite domain of elements; each tier $T_i \subseteq D$ 
	is a set of nodes. $\mathcal{T} = \{T_1,\dots,T_n\}$ is a collection of tiers such that
	$T_1 \cup T_2 \cup \dots \cup T_n = D$ and $T_i \cap T_j = \emptyset$
	for all $i \neq j$. 
	$\lhd$ is a binary successor relation defined between nodes on the same tier;
	$\alpha \subseteq D \times D$ is a set of non-crossing association relations across tiers.
\end{definition}


ASGs represent the idea that different aspects of linguistic structure
are organized on distinct tiers while progressing simultaneously over time.
Each tier encodes a particular type of information,
such as tone (pitch information), segmental (tongue shape), syllabic
(metrical) structure. Association relations link elements across tiers that are 
realized at the same temporal position.

Note that there may be multiple tiers in ASGs, provided that the 
elements on different tiers are disjoint. In most phonological 
applications, ASGs typically involve two or three tiers. In this 
work, we assume a two-tiered structure, though the definition 
remains the same in the general case. In addition, ASGs are subject 
to the No Crossing Constraint (NCC), which enforces temporal 
alignment between associations.

\begin{figure}[ht]
	\centering
	
	% ================= GRAPH =================
	\begin{minipage}[c]{0.30\textwidth}
		\vspace{0pt}
		\centering
		
		\begin{tikzpicture}[
			node/.style={draw, circle, inner sep=2pt, minimum size=18pt},
			assoc/.style={thick},
			scale=1
			]
			
			\node[node] (t1) at (0,2) {1\textsubscript{L}};
			\node[node] (t2) at (2,2) {2\textsubscript{H}};
			
			\node[node] (s1) at (0,0) {3\textsubscript{$\sigma$}};
			\node[node] (s2) at (2,0) {4\textsubscript{$\sigma$}};
			
			\draw (t1) -- node[midway, left] {$\alpha$} (s1);
			\draw (t2) -- node[midway, right] {$\alpha$} (s2);
			
			\draw (t1) -- node[midway, above] {$\lhd$} (t2);
			\draw (s1) -- node[midway, above] {$\lhd$} (s2);
			
		\end{tikzpicture}
		
	\end{minipage}
	\hfill
	% ================= MODEL =================
	\begin{minipage}[c]{0.30\textwidth}
		\vspace{0pt}
		\centering
		\footnotesize
		
		\[
		\begin{aligned}
			D &= \{1,2,3,4\} \\
			T &= \{1,2\}\\
			L &= \{1\}\\
			H &= \{2\}\\
			\sigma &= \{3,4\} \\
			\lhd &= \{
			\langle 1,2 \rangle,
			\langle 3,4 \rangle
			\} \\
			\alpha &= \{
			\langle 1,3 \rangle,
			\langle 2,4 \rangle
			\}
		\end{aligned}
		\]
		
	\end{minipage}
	\hfill
	% ================= MATRIX =================
	\begin{minipage}[c]{0.22\textwidth}
		\vspace{0pt}
		\centering
		
		\[
		\begin{array}{c|cc}
			& L & H\\
			\hline
			V_1 & 1 & 0 \\
			V_2 & 0 & 1
		\end{array}
		\]
		
	\end{minipage}
	
	\caption{
		Three representations of a two-tier autosegmental graph.
	}
	
\end{figure}
\begin{definition}[No Crossing Constraint (NCC)]
	An ASG needs to obey No Crossing Constraint. Formally,
	For all $(x,y), (x',y') \in \alpha$,
	\(
	x \lhd x' \Rightarrow (y \lhd y') \vee (y = y').
	\) and vice versa.
\end{definition}


In addition to representing model-theoretic relational structures 
with graphs, one can use \textit{adjacency matrices} which use binary values to
denote whether there is an edge between two nodes. Two nodes
are considered \textit{adjacent} if there is an edge connecting
them. 
\begin{definition}[Adjacency matrix]
	Given an ASG $\mathcal{A} = \langle D, \mathcal{T}, \alpha, \lhd \rangle$
	and two tiers $T_i, T_j \in \mathcal{T}$, the association relation $\alpha$
	can be represented as a binary matrix where 
	$M_\alpha[k,\ell] = 1$ if $(x_k, y_\ell) \in \alpha$ where $x_k \in T_i, y_\ell \in T_j$; 
	otherwise $0$.
\end{definition}

The motivation for adopting matrices is to establish a framework
corresponding to string-based learning. Concepts such as substrings,
subfactors, and superfactors are well developed and widely used in
grammar inference \refs. To incorporate ASGs into this framework,
we follow the notions of \textit{restriction} and \textit{containment} 
proposed in \citet{chandleeBufia2019}, adapting them to graph-based
representations through the notions of submatrices and matrix
containment.

\begin{definition}[Submatrix; comparable \ restriction in \citet{chandleeBufia2019}]
	\label{def.submatrix}
	\textit{
	Let matrix $A$ be an $m \times n$ matrix. A matrix $B$ is a
	\emph{submatrix} of $A$ if it is formed by taking
	$k \ (1 \leq k \leq m)$ contiguous rows and
	$v \ (1 \leq v \leq n)$ contiguous columns from $A$.
	}
\end{definition}

\begin{definition}[Containment]
\textit{
An $m \times n$ matrix $A=[a_{ij}]$ is contained
in a matrix $B$ (written $A \sqsubseteq B$) if there exists
an $m \times n$ submatrix of $B$, denoted $B'=[b'_{ij}]$,
such that for all $i \in [0,m]$ and $j \in [0,n]$,
$b'_{ij} \geq a_{ij}$.
}
\end{definition}

The containment relation provides a partial ordering over the
space of all model-theoretic structures. As illustrated by the
inverted pyramid, more general factors appear lower in the
ordering, while more specific factors appear higher positions.
This structure provides the foundation for the learning algorithm,
which determines whether a structure is contained in another.
Importantly, all factors that contain a forbidden structure, 
i.e., its superfactors, are also forbidden, while all factors 
contained in a well-formed structure, i.e., its subfactors, 
are also well formed.

\begin{figure}
\centering
\begin{tikzpicture}[
	every node/.style={inner sep=1pt, font=\small\itshape},
	edge/.style={},
	]
	\matrix (m) [matrix of nodes,
	row sep=6mm,
	column sep=10mm,
	nodes={anchor=center},
	]{
%   \node (aaa) {aaa}; &\node (llb1) {...};&   \node (llb2) {...}; & \node (llb3) {...};&\node (llb4) {...};&\node (llb5) {...};&    \node (bbb) {bbb}; \\
   &\node (aa) {aa}; & \node (ab) {ab}; && \node (ba) {ba}; & \node (bb) {bb}; \\
	 && \node (a) {a}; & & \node (b) {b}; &  \\
	 &&& \node (z) {$\epsilon$}; && \\};

	\draw[edge] (z) -- (a);

	\draw[edge] (z) -- (b);
	
	\draw[edge] (a) -- (aa) ;
	\draw[edge] (a) -- (ab);
	\draw[edge] (a) -- (ba);
	
	\draw[edge] (b) -- (bb) ;
	\draw[edge] (b) -- (ab);
	\draw[edge] (b) -- (ba);
	\draw[edge, dotted] (aa) -- ++(0,1.5em);
%	\draw[edge, dotted] (aa) -- ++(1,1.5em);
	\draw[edge, dotted] (ab) -- ++(0,1.5em);
%	\draw[edge, dotted] (ab) -- ++(1,1.5em);	
%	\draw[edge, dotted] (bb) -- ++(-1,1.5em);	
	\draw[edge, dotted] (ba) -- ++(0,1.5em);
	\draw[edge, dotted] (bb) -- ++(0,1.5em);
	\end{tikzpicture}
\caption{Strings over $\Sigma = \{a,b\}$ ordered by the containment relation (non-exhaustive)}
\label{fig.string.contain}
\end{figure}



\begin{definition}[Neighbor]
Given an ASG $\mathcal{A}=\langle D,T,V,\alpha,\lhd\rangle$, 
a node \(n\) is a neighbor of a node \(m\) iff
$m \lhd n \lor n \lhd m$.
\end{definition}

\subsection{Logical Transduction}

Using model theory \citep{strother2019diss,yolyan2025logical,jardine2017local}, 
we model a phonological transformation from an underlying 
representation (UR) to a surface representation (SR) as a 
relation between input and output structures $\mathcal{A}_i \rightarrow \mathcal{A}_o$. 
There have been work adopting the same fashion on 
string-to-string transduction \refs, but the essentially 
the logic is not restricted to strings and can operate 
over any model-theoretic representations.
In this framework, we use a set of logic formula to derive
the output signature with reference to the input signature. 
Recall from the previous section that we defined an ASG 
signature as $\mathcal{A} = \langle D, \mathcal{T}, \alpha, \lhd \rangle$.
Given an input model with this signature, the output 
signature can be defined in terms of the input structure.
\begin{table}[ht!]
\centering
\footnotesize
\begin{tabular}{lcll}
\toprule
\textbf{Output} & & \textbf{Input} & \textbf{Interpretation} \\
\midrule

$\phi_{\textsf{Domain}}$
&
$:=$
&
$\textsc{True}$
&
Preserve all input elements
\\[0.5em]


\textsc{Copy}
&
$:=$
&
1&
Define Output domain the same size as Input
\\[0.5em]


$\phi_{\textsf{license}}(x)$
&
$:=$
&
$\exists y\, \alpha(x,y)$
&
Preserve only associated elements
\\[0.5em]

$T'(x)$
&
$:=$
&
$T(x)$
&
Preserve element labels
\\[0.5em]

$\lhd'(x,y)$
&
$\approx$
&
$\lhd(x,y)$
&
Preserve linear order within-tier
\\[0.5em]

$\alpha'(x,y)$
&
$:=$
&
$\bigl(
\alpha(x,y)\land
\neg \textsf{deassoc}(x,y)
\bigr)
\lor
\textsf{reassoc}(x,y)$
&
Update association relations
\\

\bottomrule
\end{tabular}

\caption{}
\label{tab.transduction}
\end{table}

As shown in Table~\ref{tab.transduction}, the transduction 
defines an output model with the same domain as the input 
structure, while filtering out unassociated elements. All 
unary predicates \(T(x)\) in the input signature are 
preserved in the corresponding output predicates \(T'(x)\), 
meaning that an element \(x\) has property \(T\) in the 
output iff it has property \(T\) in the input. Likewise, 
the within-tier order relation is preserved. The only 
change is the association relation \(\alpha\), which is defined such 
that an association line holds in the output iff either 
(i) it is present in the input and not removed by \textsf{deassoc}, or 
(ii) it is introduced by \textsf{reassoc}. 


\[
\begin{array}{rclclcl}
\alpha'(x,y)
&:=&
\bigl(
\alpha(x,y)
&\land&
\neg \textsf{deassoc}(x,y)
\bigr)
&\lor
&
\textsf{reassoc}(x,y)
\\[0.5em]

\text{output}
&=&
\text{unchanged}
&-&
\text{delinked}
&+&
\text{newly added}
\end{array}
\]

In this way, the problem of learning association-line 
transformations can be reduced to learning the logical 
formulas for \textsf{deassoc} and \textsf{reassoc}. 
Since these two predicates are combined through disjunction, 
two parallel learning processes take place simultaneously, 
and their outputs are combined to define the final 
association relation $\alpha'$.

We demonstrate this process using a common phenomenon in tone languages
known as \textit{contour formation}. When two vowels become adjacent,
the first vowel ($V_1$) may be deleted. Since the tone originally associated
with $V_1$ cannot remain unattached, the deletion triggers delinking of
$T_1$ from $V_1$ and reassociation of $T_1$ to the following vowel $V_2$.
As a result, the surviving vowel bears both tones in their original order,
which together surface as a contour tone.

From a logical perspective, this transformation can be characterized
through the interaction of the \textsf{deassoc},
\textsf{reassoc}, and $\phi_{\textsf{license}}$ functions.
The \textsf{deassoc} operation removes the association between the
deleted vowel and its tone, indicated by the crossed-out line in
Figure~\ref{fig:hiatus}. The \textsf{reassoc} operation then links the
tone of the deleted vowel ($T_1$) to the surviving vowel ($V_2$),
indicated by the dotted line in Figure~\ref{fig:hiatus}. Consequently,
$T_1$ and $V_1$ are no longer associated and therefore are not
licensed by $\phi_{\textsf{license}}$ (shown in gray in
Figure~\ref{fig:hiatus}). Thus, $T_1, T_2$ and $V_2$ remain in the
surface representation. If we define the relevant context as
$\textsf{hiatus}(x,y) := V(x) \land V(y) \land \lhd(x,y)$,
where $x$ is the first vowel ($V_1$) and $y$ is the second vowel ($V_2$),
then the transformation can be expressed as follows:
\[
\begin{array}{ll}
	\textsf{deassoc} = \alpha(t,v) \land \textsf{delete\_V (v)}\\
	
	\textsf{reassoc} = 	\textsf{survive\_V} \land\exists v\,(\alpha(t,v)\land \textsf{delete\_V})
\end{array}
\]

\begin{figure}[ht!]
	\centering
	
	\begin{minipage}[t]{0.48\textwidth}
		\centering
		
		\begin{tabular}{c c c}
			
			\textbf{UR} & $\Rightarrow$ & \textbf{SR} \\[0.5em]
			
			\begin{tikzpicture}[scale=0.7, every node/.style={inner sep=2pt}]
				\node [fill=gray!20](v1) at (0,0) {V$_1$};
				\node (v2) at (1.4,0) {V$_2$};
				
				\node (t1) at (0,1.3) {T$_1$};
				\node (t2) at (1.4,1.3) {T$_2$};
				
				\draw (v1)--(t1);
				\draw (v2)--(t2);
			\end{tikzpicture}
			
			&&
			
			\begin{tikzpicture}[scale=0.7, every node/.style={inner sep=2pt}]
				\node [fill=gray!20](v1) at (0,0) {V$_1$};
				\node (v2) at (1.4,0) {V$_2$};
				
				\node (t1) at (0,1.3) {T$_1$};
				\node (t2) at (1.4,1.3) {T$_2$};
				
				\draw (v1)--(t1) node[midway] {=};
				\draw (v2)--(t2);
				
				\draw[dotted, thick] (v2)--(t1);
			\end{tikzpicture}
			
		\end{tabular}
		
	\end{minipage}
	\hfill
	\begin{minipage}[t]{0.48\textwidth}
		\centering
		
		\begin{tabular}{c c c}
			
			\textbf{UR} & $\Rightarrow$ & \textbf{SR} \\[0.5em]
			
			$
			\begin{array}{c|cc}
				& T_1 & T_2\\
				\hline
				V_1 & \cellcolor{gray!75}1 & 0 \\
				V_2 & \cellcolor{gray!15}0 & 1 \\
			\end{array}
			$
			
			&&
			
			$
			\begin{array}{c|cc}
				& T_1 & T_2\\
				\hline
				V_1 & \cellcolor{gray!15}0 & 0 \\
				V_2 & \cellcolor{gray!75}1 & 1 \\
			\end{array}
			$
			
		\end{tabular}
		
	\end{minipage}
	
	\caption{
%		Vowel hiatus resolution and the corresponding adjacency matrices.
	}
	\label{fig:hiatus}
	
\end{figure}

The \textsf{deassoc} and \textsf{reassoc} operations are
equivalently reflected in the adjacency matrices.
Figure~\refs shows the input and output
matrices corresponding to the structures in
Figure~\ref{fig:hiatus}. As can be seen, the value
corresponding to the association between $T_1$ and $V_1$
changes from $1$ in the input to $0$ in the output,
reflecting delinking. Conversely, the value corresponding
to $T_1$ and $V_2$ changes from $0$ to $1$,
reflecting reassociation. In this way, the autosegmental
graph, the adjacency matrix, and the logical expressions
provide equivalent description of the same
transformation. Formally, \textsf{deassoc} and \textsf{reassoc}
can also be identified by  $ \textsf{deassoc}
= \{(i,j) \mid \mathrm{UR}_{ij}=1 \land \mathrm{SR}_{ij}=0\}$
and $\textsf{reassoc} =
\{(i,j) \mid \mathrm{UR}_{ij}=0 \land \mathrm{SR}_{ij}=1\}$.
\subsection{Association Learning}
 This section follows the framework of using logical transductions
to learn phonological mappings in \citet{yolyan2025logical},
which leverages the entailment relation
in a partially ordered space and updates hypotheses about
contextual environments. One important aspect of the framework
is the use of two sets of $k$-factors: $V(X)$ and
$\hat{V}(X)$x, which contain the relevant $k$-factors
associated with the target transformation. 
$V(X)$ contains the \textit{positive evidence}, namely all
$k$-factors in which the target transformation takes place
from the input to the output. In contrast,
$\hat{V}(X)$ contains all \textit{negative evidence} where
the target transformation does not take place, that is,
where the relevant property remains unchanged between
the input and output structures. 

\citet{yolyan2025logical} demonstrates how to learn a 
string-based input-output mapping that can be written as
 $A \rightarrow B \,/\, C \_ D$, where the goal
of learning is to infer the contextual environment
under which the transformation applies. She defines

\begin{definition}
For every $X \in \{\Phi_F,\varphi_F\}_{F \in \mathcal{F}}$,
the corresponding hypothesis space is defined as
\[
\mathcal{H}(X) :=
\bigcup_{x \in V(X)} \downarrow x
\;-\;
\bigcup_{x \in \hat{V}(X)} \downarrow x .
\]
\end{definition}

In autosegmental phonology, such
processes are not always easy to express linearly, since the relevant
structure involves relations rather than simple strings. However, the
learning question remains the same: what environment triggers the
change from \(1\) to \(0\), and what environment triggers the change
from \(0\) to \(1\)?







With only a single UR--SR pair, there are infinitely many possible
hypotheses. For example, one might hypothesize that all \(H\) tones
delete their associations, or that all vowels do. However, once more
data points are observed, counterexamples eliminate incorrect
generalizations.

The learning task, therefore, is to determine the surrounding
environment of a particular cell value—either \(1\) or \(0\)—that
triggers the change.

\[
\begin{array}{ccc}
	? & ? & ? \\
	? & \colorbox{pink}{0} & ? \\
	? & ? & ?
\end{array}
\quad \rightarrow \quad
\begin{array}{ccc}
	? & ? & ? \\
	? & \colorbox{pink}{1} & ? \\
	? & ? & ?
\end{array}
\]

\begin{algorithm}[th!]
\small
\KwIn{
	$
	D = \{(I_1,O_1),\dots,(I_n,O_n)\}
	$
}
\KwOut{
	A hypothesis $H$
}

$H,V^+,V^-,P \gets \emptyset$\;

$Q \gets \{\alpha(x,y)\}$\;

\ForEach{$(I,O)\in D$}{
	
	$V^+ \gets V^+ \cup (\alpha_I \setminus \alpha_O)$\;
	
	$V^- \gets V^- \cup (\alpha_I \cap \alpha_O)$\;
}

\While{$Q \neq \emptyset$}{
	
	$c \gets Q.\mathrm{Dequeue}()$\;
	
	\If{$\forall g\in V^+,\; c\sqsubseteq g$}{
		
		\If{$\forall g\in V^-,\; c\not\sqsubseteq g$}{
			
			$H \gets H \cup \{c\}$\;
		}
		
		\Else{
			
			\ForEach{$s\in \textsc{NextSupFact}(c)$}{
				
				\If{$\forall p\in P,\; s\not\sqsubseteq p$}{
					
					$Q.\mathrm{Enqueue}(s)$\;
				}
			}
		}
	}
	
	\Else{
		
		$P \gets P \cup \{c\}$\;
	}
}

\Return{$H$}\;

\end{algorithm}

\section{Case Study}
\subsection{Contour Formation}
We begin with an idealized assumption: the tonal process involves no
change to the segmental string, and the only modifications between the
underlying representation (UR) and the surface representation (SR) are
changes in association lines, namely \textsc{deassociation} and
\textsc{reassociation}. Although many phonological processes also
involve segmental deletion, insertion, or feature changes, we restrict
the current discussion to association-line modifications in order to
isolate the learning problem.

The first step is to identify exactly what structural modification has
taken place between the UR and the SR. Since the segmental and tonal
tiers remain constant, the comparison is made directly over the
association relation matrix.

For each cell in the association matrix:


That is, if an association present in the UR disappears in the SR, the
operation is classified as deassociation; if a new association appears
in the SR that was absent in the UR, the operation is classified as
reassociation.

In Estoko, both operations are attested: one association line is
removed while another is added. Therefore, the learner must infer both
the deassociation environment and the reassociation environment
separately. (maybe not maybe they can start together)
\paragraph{Hypothesis space for deassociation.}

Given a set of observed deassociation sites, the learner constructs
candidate hypotheses describing the structural environments in which an
association changes from $1$ to $0$. Each hypothesis corresponds to a
local structural condition under which $\alpha(x,y)$ is removed.

The search proceeds from simple to more complex descriptions. At each
step, the learner considers increasingly rich local configurations
(factors) until a hypothesis is found that captures all and only the
observed deassociation sites.

The most general hypothesis removes all associations:
\[
\text{Hypothesis 1: } 1 \rightarrow 0
\]
This hypothesis is immediately rejected if there exist associations in
the data that remain unchanged.


\section{Conclusion} This paper has argued that vowel-deletion–induced 
tonal processes can be uniformly captured as logical transductions 
over autosegmental representations (ARs). The observed typological 
variation follows from modifying a single relation, $\alpha$, through 
the interaction of \textsf{deassoc} and \textsf{reassoc}. 

A limitation of the current analysis is that the transduction relies 
on quantification to capture across-tier dependencies. An important 
direction for future work is to determine whether these processes can 
be simplified in a quantifier-free (QF) process.
This work also extends naturally to learning over ARs. Recent work 
by \citet{yolyan2025logical} demonstrates that logical transductions 
are learnable by using a bottom-up factorial inference algorithms\citep{chandleeBufia2019}. This opens up more discussions about how linguistically meaningful 
representations can facilitate grammar inference.


\bibliographystyle{apalike}
\bibliography{pmlr-sample}

\end{document}